%% ma: L12-B08-0167
%% lop: 12
%% chuong: 2
%% bai: 8
%% dang: 12.8.30
%% loai: TN4
%% muc: VD
%% dap_an: "B"
%% nguon: "CĐ2.4.pdf giáo viên gửi 27/10/2026 (M07, Lê Minh Khôi – Nguyễn Anh Khôi, Chương II Chủ đề 4 Các bài toán về tâm tỉ cự), Câu 17"
%% kien_thuc: "Bài 8 (khoảng cách hai điểm, biểu thức tọa độ), Bài 6 (bình phương vô hướng, tâm tỉ cự)"
%% trang_thai: chinh-thuc
\begin{ex}
Trong không gian $Oxyz$, cho ba điểm $A(1;1;1)$, $B(-1;2;1)$, $C(3;6;-5)$. Biết rằng $M(a;b;c)$ là điểm thuộc mặt phẳng $(Oxy)$ thỏa mãn $MA^2+MB^2+MC^2$ đạt giá trị nhỏ nhất. Tính giá trị của biểu thức $J=a+b+c$.
\choice{$J=3$}{\True $J=4$}{$J=5$}{$J=6$}
\loigiai{Gọi $G(1;3;-1)$ là trọng tâm tam giác $ABC$: $MA^2+MB^2+MC^2=3MG^2+GA^2+GB^2+GC^2$.

Nhỏ nhất khi $M$ là hình chiếu của $G$ lên $(Oxy)$: $M(1;3;0)$, $J=4$.}
\end{ex}
